\documentclass{article}
\usepackage[utf8]{inputenc}
\usepackage[T1]{fontenc}
\usepackage[margin=2cm]{geometry}
\usepackage{hyperref}
\sloppy
\begin{document}
\section*{Historical source}
Encyclopaedia Britannica scan, source row 851. Parse 3.5 extracted the text. Super 3.5 VL corrected it from an enlarged crop of each block. It can still contain recognition errors.
Page headers, footers, and page numbers are removed.

\subsection*{OPTICS (E0071)}
OPTICS.

Part II.

Of
Reflection.

tre is C, and focus of parallel rays is at P; and let the imaginary focus of the incident rays be at F, a point between P and B; and let DR be the reflected ray. From C and R draw the lines CH, RH, the one passing through D, the other parallel to FE. Then will the angle RHD be equal to HDE the angle of incidence, as alternate to it; and therefore equal to HDR, the angle of reflection: wherefore the triangle HDR is isosceles, and consequently DR is equal to RH. Now the lines FD and RH being parallel, the triangles FDC and RHC are similar; and therefore RH, or RD, is to FD as CR to CF: but BD vanishing, RD and FD coincide with RB and FB, wherefore RB is to FB as CR to CF; that is, the distance of the focus from the surface is to the distance of the imaginary focus of the incident rays, as the distance of the focus from the centre is to the distance of the imaginary focus of the incident rays from the same. Q. E. D.

Case 6. Of diverging rays falling upon a concave surface, and proceeding from a point between the focus of parallel rays and the surface.

Demon. Let FD and FB represent two diverging rays flowing from the point F as a radiant, which was the imaginary focus of the incident rays in the foregoing case. Then as ED was in that case reflected into DR, FD will be reflected into DK, (for the reason mentioned in Case the second), so that the reflected ray will proceed as from the point R: but it was demonstrated in the case immediately foregoing, that RB is to FB as CR to CF; that is, the distance of the focus from the surface is to that of the radiant from the same, as the distance of the former from the centre is to that of the latter from the same. Q. E. D.

Case 7. Of converging rays falling upon a convex surface, and tending towards a point beyond the centre.

Dem. Let AB, ED, (fig. 8.) be the incident rays tending to F, a point beyond the centre C, and let DK be the reflected ray of the incident one ED. Then because the incident ray ED tends to a point beyond the centre, the reflected ray DK will proceed as from one on the contrary side, suppose R; as explained above under Prop. VII. Through D draw the perpendicular CD, and produce it to H. Then will EDH and HDK be the angles of incidence and reflection; which being equal, their vertical ones CDF and CDR will be fo too: consequently the vertex of the triangle FDR is bisected by the line CD: wherefore, RD is to DF, or (3 Elem. 6.) BD vanishing, RB is to BF as RC to CF; that is, the distance of the focus of the reflected rays is to that of the imaginary focus of the incident rays, as the distance of the former from the centre is to the distance of the latter from the fame. Q. E. D.

CASE 8. Of diverging rays falling upon a concave surface, and proceeding from a point beyond the centre.

Dem. Let FB, FD, be the incident rays having their radiant in F, the imaginary focus of the incident rays in the foregoing case. Then as ED was in that case reflected into DK, FD will now be reflected into DR; fo that R will be the focus of the reflected rays. But it was demonstrated in the foregoing case, that RB is to FB as RC to CF; that is, the distance of

the focus of the reflected rays from the surface is to the distance of the radiant from the same, as the distance of the focus of the reflected rays from the centre is to the distance of the radiant from thence. Q. E. D.

The two remaining cafes may be confirmed as the converse of thote under Prop. II. (p. 548-), because the incident rays in there are the reflected ones in them; or they may be demothered in the fame manner with the forgoing, as follows.

Case 9. Converging rays falling upon a convex surface, and tending to the surface of parallel rays, become parallel after reflection.

Dem. Let ED, RB, (fig. 7.) represent two converging rays incident on the convex surface BD, and tending towards F, which we will now suppose to be the focus of parallel rays; and let DR be the reflected ray, and C the centre of convexity of the reflecting surface. Through C draw the line CD, and produce it to H, drawing RH parallel to ED produced to F. Now it has been demonstrated (Case 5.) where the incident rays are supposed to tend to the point F, that RB is to FB as RC to CF; but F in this Case being supposed to be the focus of parallel rays, it is the middle point between C and B (by proposition 2d), and therefore FB and FC are equal; and consequently the two other terms in the proportion, viz. RB and RC, must be fo too; which can only be upon a supposition that R is at an infinite distance from B; that is, that the reflected rays BR and DR be parallel. Q. E. D.

CASE 10. Diverging rays falling upon a concave surface, and proceeding from the focus of parallel rays, become parallel after reflection.

Demon. Let RD, RB (fig. 8.) be two diverging rays incident upon the concave surface BD, as supposed in Case 4.; where it was demonstrated that FB is to RB as CF to CR. But in the present case RB and CR are equal, because R is supposed to be the focus of parallel rays; therefore FB and FC are fo too. Which cannot be unless F be taken at an infinite distance from B; that is, unless the reflected rays BF and DF be parallel. Q. E. D.

It is here observable, that in the case of diverging rays falling upon a convex surface, (fee fig. 7) the farther the point D is taken from B, the nearer the point F, the focus of the reflected rays, approaches to B, while the radiant R remains the fame. For it is evident from the curvature of a circle, that the point D (fig. 9) may be taken so far from B, that the reflected ray DE shall proceed as from F, G, H, or even from B, or from any point between B and R; and the further it is taken from B, the faster the point from which it proceeds, approaches towards R: as will easily appear if we draw several incident rays with their respective reflected ones, in such manner that the angles of reflection may be all equal to their respective angles of incidence, as is done in the figure. The like is applicable to any of the other cases of diverging or converging rays incident upon a spherical surface. This is the reason, that when rays are considered as reflected from a spherical surface, the distance of the oblique rays from the perpendicular one is taken so small, that it may be supposed to vanish.

From


\section*{Source record}
\begin{description}
\item[Scan ID] \texttt{britannica:309f765b1ecc83e0a7f7762e44b4a927a7555443:train:851}
\item[Image file] \texttt{page-0851-6e2a0b999324.jpg}
\item[Image SHA-256] \texttt{6e2a0b9993242da4f83259432110d6929e728a4c0e58a9e79a16ff80a6fd1861}
\item[Dataset] NationalLibraryOfScotland/encyclopaedia\_britannica\_illustrated at revision \texttt{309f765b1ecc83e0a7f7762e44b4a927a7555443} (CC0-1.0)
\item[Extraction] model \texttt{nvidia/nemotron-3.5-parse}, raw reply \texttt{extract/raw/page-0851-6e2a0b999324.json}
\item[Raw reply SHA-256] \texttt{7101b44c603aa59acabd133ee1fba9529a293ccd4087adfdb4f44bff255ef705}
\item[Correction] model \texttt{nvidia/Nemotron-3.5-Super-VL}, block crops, raw replies \texttt{correct/raw/crop\_nothink/page-0851-6e2a0b999324-b*.json}; super\_vl: 18
\item[Running header] none
\end{description}
\end{document}
